(1) 小振幅张紧薄膜的横位移 \(w(r,\phi,t)\) 满足 \(\rho_s w_{tt}=T\nabla_{\!2}^2w\)，其中极坐标Laplace算子为 \(\nabla_{\!2}^2=\partial_r^2+(1/r)\partial_r+(1/r^2)\partial_\phi^2\)。边界 \(w(a,\phi,t)=0\)，波速 \(c=\sqrt{T/\rho_s}\)。
(2) 对于与 \(\phi\) 无关的轴对称振动，令 \(w=R(r)e^{i\omega t}\) 得 \(R''+R'/r+k^2R=0\)，\(k=\omega/c\)。它是零阶Bessel方程，一般解为 \(AJ_0(kr)+BY_0(kr)\)；第二种函数在 \(r=0\) 对数发散，正则性要求 \(B=0\)。
(3) 固定圆边给 \(J_0(ka)=0\)，记 \(j_{0,n}\) 为 \(J_0\) 的第 \(n\) 个正零点，故 \(\omega_n=cj_{0,n}/a\)。最低轴对称频率为 \(\omega_1\simeq2.404826\,c/a\)，随着膜半径增加按 \(a^{-1}\) 降低。
(4) 振动能量为 \(E=\frac12\int_0^{2\pi}\int_0^a[\rho_s w_t^2+T|\nabla_{\!2}w|^2]rdrd\phi\)。积分分部并用固定边缘 \(w_t(a)=0\)，可得 \(\dot E=0\)；若角向依赖为 \(e^{im\phi}\)，径向方程出现 \(m^2/r^2\)，其正则解是 \(J_m(kr)\)。
